\documentclass{article}
\usepackage{amsmath,amssymb,amsthm}
\usepackage{graphicx}
\usepackage{hyperref}
\usepackage[margin=1in]{geometry}

\newtheorem{proposition}{Proposition}
\newtheorem{claim}{Claim}
\newtheorem{corollary}{Corollary}
\newcommand{\btheta}{\boldsymbol{\theta}}
\newcommand{\bSigma}{\boldsymbol{\Sigma}}
\newcommand{\bH}{\mathbf{H}}
\newcommand{\R}{\mathbb{R}}
\newcommand{\E}{\mathbb{E}}
\newcommand{\Var}{\mathrm{Var}}

\title{A Theoretical Model of Label-Noise-Driven Training Trajectory Divergence \\
and Linear Mode Connectivity Barrier Dynamics}
\author{}
\date{}

\begin{document}
\maketitle

\section{Overview}

We provide a formal optimization-dynamics model that explains why
medical-domain fine-tuning produces monotonically growing linear-mode-connectivity
(LMC) barriers while code-domain fine-tuning exhibits an inverted-U barrier
profile that declines during the low-learning-rate ``putt'' phase.  The key
mechanism is \emph{label-noise-induced gradient misalignment}, which prevents
the low-LR convergence back to the pre-trained basin that reconnects code-domain trajectories.

\section{Empirical Phenomena to Explain}

\begin{enumerate}
\item \textbf{Code domain}: barrier peaks at 0.053 (step $\sim$200, $\Delta W \approx 1.5\%$),
      then declines to 0.033 at step 400.  Inverted-U shape.
\item \textbf{Medical domain}: barrier grows monotonically $0.173 \to 0.218$,
      plateaus without decline.
\item \textbf{Gaussian noise} applied to weights produces negligible barrier
      ($\max 0.014$) even at $\Delta W = 8\%$, while training produces 0.053 at
      $\Delta W = 1.5\%$ --- a factor of 9--20$\times$ per unit $\Delta W$.
\item \textbf{Early layers (0--7)} account for 75\% of total barrier; late
      layers (16--23) have negligible barrier.
\item \textbf{Seed-pair asymmetry}: seed $s_4$ paired with $s_1$ produces
      barrier 1.213, but with $s_2$ produces 0.080.
\item \textbf{Cross-domain barriers} are consistent ($\sim$0.05) regardless
      of domain pair, while within-domain barriers span 0.048 (code) to
      0.147 (medical).
\item \textbf{At intermediate $\Delta W \sim 5\%$}, catastrophic barriers
      (1.043) coexist with best self-domain performance.
\end{enumerate}

\section{SGD as a Stochastic Differential Equation}

\subsection{Setup}

Consider a neural network $f_\btheta: \mathcal{X} \to \R^K$ parameterized by
$\btheta \in \R^d$ ($d \approx 1.4 \times 10^9$ for Pythia-1.4B).  We have
a training set $\mathcal{D} = \{(\mathbf{x}_i, y_i)\}_{i=1}^N$.  The
population loss is $L(\btheta) = \E_{(\mathbf{x},y) \sim P}[\ell(f_\btheta(\mathbf{x}), y)]$
and the empirical loss is $\hat{L}(\btheta) = \frac{1}{N}\sum_i \ell(f_\btheta(\mathbf{x}_i), y_i)$.

SGD with mini-batch $\mathcal{B}_t$ of size $B=128$ and learning rate $\alpha_t$:
\begin{equation}
\btheta_{t+1} = \btheta_t - \alpha_t \hat{\mathbf{g}}_t, \qquad
\hat{\mathbf{g}}_t = \frac{1}{B}\sum_{i \in \mathcal{B}_t} \nabla_\btheta \ell(f_{\btheta_t}(\mathbf{x}_i), y_i).
\end{equation}

For small $\alpha_t$, this is well approximated by the continuous-time SDE (Li et al., 2017;
Mandt et al., 2017):
\begin{equation}
d\btheta = -\nabla L(\btheta)\,dt + \sqrt{\alpha_t\,\bSigma(\btheta)}\,d\mathbf{W}_t,
\label{eq:sde}
\end{equation}
where $\bSigma(\btheta) = \frac{1}{B}\mathrm{Cov}_{\mathcal{D}}[\nabla_\btheta \ell(f_\btheta(\mathbf{x}), y)]$
is the gradient noise covariance matrix, and $\mathbf{W}_t$ is $d$-dimensional Brownian motion.

\subsection{The Drive--Putt LR Schedule}

The paper uses a cosine schedule with ``drive--putt'' structure:
\begin{equation}
\alpha_t = \begin{cases}
\alpha_{\text{drive}} \approx 1 \times 10^{-4}, & t \in [0, 200] \\
\alpha_{\text{putt}} \approx 1 \times 10^{-5},  & t \in [200, 400]
\end{cases}
\end{equation}

The qualitative behavior differs sharply between phases:
\begin{itemize}
\item \textbf{Drive phase} ($\alpha_t$ large): The diffusion term
      $\sqrt{\alpha_t \bSigma}$ dominates.  SGD explores the loss landscape
      broadly.  Two independent seeds undergo distinct random walks, producing
      weight divergence $\Delta W(t)$.
\item \textbf{Putt phase} ($\alpha_t$ small): The drift term $-\nabla L$
      dominates \emph{provided the gradient signal is clean}.  SGD converges
      toward the nearest local minimum.  If two seeds land in overlapping
      basins of attraction during the drive phase, they converge to nearby
      minima during the putt phase, and the LMC barrier \emph{decreases}.
\end{itemize}

\section{Label Noise and Gradient Misalignment}

\subsection{PRM Verification Data: Sources of Label Noise}

In PRM (Process Reward Model) verification, each training example is a
reasoning step annotated as $\{\text{correct}, \text{incorrect}\}$.

\begin{itemize}
\item \textbf{Code domain}: Correctness is deterministically verifiable by
      execution.  Label noise is negligible: $\Pr(y_{\text{true}} \neq y_{\text{annotated}} | \mathbf{x}) \approx 0$.
\item \textbf{Medical domain}: Step correctness is inherently ambiguous.
      Clinical reasoning admits multiple valid pathways; annotator
      disagreement is common.  We model this as:
      \begin{equation}
      \Pr(y_{\text{true}} = \text{correct} \mid \mathbf{x}) = p(\mathbf{x}),
      \end{equation}
      where $p(\mathbf{x}) \in (0.4, 0.6)$ for a substantial fraction of
      examples.  This is \emph{maximum-entropy label noise}: the observed
      label $y$ carries minimal information about the true function.
\end{itemize}

\subsection{Gradient Noise Covariance Under Label Noise}

Let $g(\btheta; \mathbf{x}, y) = \nabla_\btheta \ell(f_\btheta(\mathbf{x}), y)$.
Under label noise with flip probability $\varepsilon(\mathbf{x}) =
\min(p(\mathbf{x}), 1-p(\mathbf{x}))$:

\begin{proposition}[Label-noise inflation of gradient covariance]
\label{prop:noise}
For binary classification with cross-entropy loss, the gradient noise
covariance decomposes as:
\begin{equation}
\bSigma(\btheta) = \underbrace{\bSigma_{\text{clean}}(\btheta)}_{\text{sampling noise}} \;+\;
\underbrace{\bSigma_{\text{label}}(\btheta)}_{\text{label noise}},
\end{equation}
where
\begin{equation}
\bSigma_{\text{label}}(\btheta) \approx \frac{1}{B}\,\E_{\mathbf{x}}\Big[
\varepsilon(\mathbf{x})(1 - \varepsilon(\mathbf{x}))\,
\delta g(\mathbf{x}) \, \delta g(\mathbf{x})^\top \Big],
\end{equation}
with $\delta g(\mathbf{x}) = g(\btheta; \mathbf{x}, \text{correct}) - g(\btheta; \mathbf{x}, \text{incorrect})$
being the gradient difference between the two label assignments.

The key factor $\varepsilon(1-\varepsilon)$ is maximized at $\varepsilon = 0.5$,
giving:
\begin{equation}
\frac{\|\bSigma_{\text{label}}^{\text{med}}\|}{\|\bSigma_{\text{label}}^{\text{code}}\|}
\approx \frac{0.5 \times 0.5}{0.05 \times 0.95} \approx 5.3\times.
\end{equation}
\end{proposition}

\subsection{Effective Diffusion During Drive Phase}

The weight divergence between two seeds $A, B$ after the drive phase
($T_{\text{drive}} = 200$ steps) is approximately:
\begin{equation}
\E[\|\btheta_T^A - \btheta_T^B\|^2] \approx
2\alpha_{\text{drive}} T_{\text{drive}} \, \mathrm{tr}(\bSigma(\btheta_0)).
\label{eq:divergence}
\end{equation}

\begin{claim}
The medical gradient noise's 5--10$\times$ inflation of $\mathrm{tr}(\bSigma)$
directly produces larger $\Delta W$ at the end of the drive phase, and
\emph{qualitatively different} directions of divergence (less aligned with
the true loss gradient, more random).
\end{claim}

\section{Convergence vs.\ Non-Convergence During Putt Phase}

\subsection{The Stationary Distribution}

Near a local minimum $\btheta^*$, the loss is approximately quadratic:
$L(\btheta) \approx L(\btheta^*) + \frac{1}{2}(\btheta - \btheta^*)^\top \bH (\btheta - \btheta^*)$.
The SDE~\eqref{eq:sde} becomes an Ornstein-Uhlenbeck process:
\begin{equation}
d\btheta = -\bH(\btheta - \btheta^*)\,dt + \sqrt{\alpha_{\text{putt}}\,\bSigma(\btheta^*)}\,d\mathbf{W}_t.
\end{equation}

The stationary distribution is $\mathcal{N}(\btheta^*, \alpha_{\text{putt}} \bH^{-1} \bSigma \bH^{-1})$.
Two key metrics:

\begin{enumerate}
\item \textbf{Convergence time}: $\tau_{\text{conv}} \sim 1/(\alpha_{\text{putt}} \lambda_{\min}(\bH))$.
      With $\alpha_{\text{putt}} = 10^{-5}$ and typical $\lambda_{\min}(\bH) \sim 10^{-2}$,
      $\tau_{\text{conv}} \sim 10^7$ gradient steps.  The putt phase (200 steps)
      is far too short for full convergence.

\item \textbf{Stationary spread}: $\sigma_{\text{stat}} \sim
      \sqrt{\alpha_{\text{putt}} \, \|\bH^{-1} \bSigma \bH^{-1}\|}$.
      Larger $\bSigma$ (medical) $\implies$ wider stationary distribution
      $\implies$ models from different seeds are more likely to occupy
      \emph{different} minima.
\end{enumerate}

\subsection{The Peclet Number: Drift vs.\ Diffusion}

Define the local Peclet number comparing deterministic drift to stochastic diffusion:
\begin{equation}
\mathrm{Pe}(\btheta) = \frac{\|\nabla L(\btheta)\|}{\sqrt{\alpha_{\text{putt}} \, \mathrm{tr}(\bSigma(\btheta))}}.
\end{equation}

\begin{itemize}
\item $\mathrm{Pe} \gg 1$: Drift dominates.  SGD moves deterministically
      toward the nearest minimum.  Multiple seeds in the same basin converge
      to nearby points.  Barrier \emph{decreases}.
\item $\mathrm{Pe} \lesssim 1$: Diffusion competes with drift.  SGD wanders
      without converging to a specific minimum.  Different seeds explore
      different regions.  Barrier \emph{does not decrease}.
\end{itemize}

\begin{proposition}[Peclet ratio across domains]
Under the label-noise model of Proposition~\ref{prop:noise}:
\begin{equation}
\frac{\mathrm{Pe}_{\text{code}}}{\mathrm{Pe}_{\text{medical}}}
\approx \sqrt{\frac{\mathrm{tr}(\bSigma_{\text{med}})}{\mathrm{tr}(\bSigma_{\text{code}})}} \approx 2.3.
\end{equation}
At $\alpha_{\text{putt}} = 10^{-5}$, code models satisfy $\mathrm{Pe} \gg 1$
in most of the relevant region; medical models are marginal ($\mathrm{Pe} \sim 1$)
in large portions of the loss landscape.
\end{proposition}

\subsection{Putt-Phase Dynamics: The Full Picture}

During the putt phase (steps 200--400, $\sim$200 gradient updates):

\noindent\textbf{Code models} ($\mathrm{Pe} \gg 1$, small $\bSigma$):
\begin{itemize}
\item Drift dominates.  $\btheta_t$ moves deterministically toward
      $\btheta^*_{\text{code}}$, a minimum in or near the pre-trained basin.
\item Two seeds $A, B$ both converge to nearby minima $\btheta^*_A, \btheta^*_B$
      in the same broad basin.
\item The linear interpolation barrier $B(\btheta^A_{400}, \btheta^B_{400})$
      decreases relative to $B(\btheta^A_{200}, \btheta^B_{200})$.
\item This produces the \textbf{inverted-U}: barrier rises during drive
      (seeds diverge), then falls during putt (seeds converge to same basin).
\item Quantitatively: barrier drops from 0.053 $\to$ 0.033, a 38\%
      reduction in 200 steps.
\end{itemize}

\noindent\textbf{Medical models} ($\mathrm{Pe} \sim 1$, large $\bSigma$):
\begin{itemize}
\item Diffusion is comparable to drift.  $\btheta_t$ does not converge
      cleanly to any particular minimum.
\item Two seeds $A, B$ wander quasi-randomly, landing in \emph{different}
      local minima whose distance grows over time.
\item The barrier $B(\btheta^A_t, \btheta^B_t)$ grows monotonically because:
      \begin{enumerate}
      \item Each seed's path is dominated by different noise realizations.
      \item The minima they approach are in different, disconnected basins.
      \item Label-noise-induced gradient misalignment means even the
            \emph{expected} gradient points in different directions for
            different seeds at the same $\btheta$ (the noise is
            \emph{state-dependent} as well as seed-dependent).
      \end{enumerate}
\item Quantitatively: barrier grows from 0.173 $\to$ 0.218, a 26\%
      \emph{increase} over the same 200 steps.
\end{itemize}

\section{Why Gaussian Noise Produces Negligible Barriers}

The empirical finding that Gaussian weight perturbation ($\max$ barrier 0.014
at $\Delta W = 8\%$) is far weaker than training-induced barriers (0.053 at
$\Delta W = 1.5\%$) has a clear theoretical explanation:

\begin{claim}[Isotropic vs.\ structured perturbations]
Gaussian noise $\boldsymbol{\epsilon} \sim \mathcal{N}(0, \sigma^2 \mathbf{I}_d)$
is \textbf{isotropic} in parameter space.  The network function $f_\btheta$
is highly insensitive to isotropic perturbations because:
\begin{enumerate}
\item Most directions in parameter space are ``flat'' (near-zero Hessian
      eigenvalues) and do not affect the function.
\item The network can internally compensate: a random perturbation in
      layer $\ell$ is partially canceled by adjustments in layer $\ell+1$
      during the forward pass.
\item The effective functional change $\|f_{\btheta + \boldsymbol{\epsilon}} - f_\btheta\|$
      scales as $\mathcal{O}(\|\boldsymbol{\epsilon}\|/\sqrt{d})$ due to
      concentration of measure in high dimension ($d \sim 10^9$).
\end{enumerate}
\end{claim}

In contrast, SGD noise is \textbf{structured}:
\begin{enumerate}
\item It is concentrated in early layers (layers 0--7), which control
      fundamental representations.  Small perturbations here propagate
      nonlinearly through all subsequent layers.
\item It is aligned with the gradient, which is the direction of
      \emph{maximal} functional change per unit parameter change.
\item It is correlated across parameters within each layer, producing
      coherent rather than self-canceling changes.
\end{enumerate}

A simple scaling argument: if Gaussian noise of magnitude $\delta$ changes
the function by $\sim \delta/\sqrt{d_{\text{eff}}}$ (where $d_{\text{eff}}$
is the effective dimension of the parameter space), while SGD noise of the
same magnitude changes the function by $\sim \delta \cdot \|\nabla_\btheta f\|$
(which is $\mathcal{O}(1)$ in the gradient direction), then the ratio of
functional impact is $\sim \sqrt{d_{\text{eff}}} \sim 10^4$, explaining why
training-induced barriers are $\sim$10--100$\times$ larger per unit $\Delta W$.

\section{Early-Layer Concentration of Barriers}

The empirical finding that layers 0--7 account for 75\% of the total barrier
has a natural explanation in our framework.

\begin{proposition}[Layer-wise gradient noise structure]
For transformer architectures, the gradient norm $\|\nabla_{\btheta_\ell} L\|$
and the gradient noise covariance $\mathrm{tr}(\bSigma_\ell)$ are both
concentrated in early layers during fine-tuning.  This is because:
\begin{enumerate}
\item Early layers learn input-formatting representations that must adapt
      to the new domain (reasoning traces vs.\ pre-training text).
\item Changes in early layers propagate multiplicatively through the
      self-attention and FFN blocks of all subsequent layers.
\item The effective ``gain'' from layer $\ell$ to the output is
      $\sim \prod_{k=\ell}^{L} \|\mathbf{W}_k\|$, which is exponentially
      larger for small $\ell$.
\end{enumerate}
\end{proposition}

Thus, the same amount of parameter noise in layer 0 produces exponentially
larger functional change than in layer 23, and the LMC barrier (which measures
functional mismatch along the linear interpolation) is correspondingly
concentrated in early layers.

\section{Seed-Pair Compatibility: A Basin-Membership Model}

The stark asymmetry in seed-pair barriers ($s_4$--$s_1$: 1.213 vs.\
$s_4$--$s_2$: 0.080) can be modeled as a \textbf{basin-membership} problem.

During the drive phase, each seed's trajectory is a random walk with
state-dependent diffusion.  Whether two seeds land in the same basin of
attraction of the putt-phase dynamics depends on:
\begin{enumerate}
\item The \textbf{initial noise realization} in the first few steps,
      which sets the ``direction'' of the random walk.
\item The \textbf{loss landscape geometry}: basins have finite width
      $\sim R_{\text{basin}}$ in parameter space.
\item If the distance between seeds at the drive--putt transition
      exceeds $R_{\text{basin}}$, they fall into different basins and
      the putt phase cannot reconnect them.
\end{enumerate}

\begin{claim}
For code models, $R_{\text{basin}}$ is large (clean gradient signal $\to$
broad, well-defined basin).  Most seed pairs land in the same basin
$\to$ low within-domain barrier (0.048 on average).  For medical models,
$R_{\text{basin}}$ is small (noisy gradient $\to$ fragmented landscape with
many sharp local minima).  Seed pairs routinely land in different basins
$\to$ high within-domain barrier (0.147 on average).  The asymmetry
($s_4$ with $s_1$ vs.\ $s_2$) reflects whether the pair happened, by
chance, to land in the same basin or different basins.
\end{claim}

\section{Cross-Domain Barriers: The Pre-Trained Basin as a Universal Anchor}

The consistent cross-domain barrier ($\sim$0.05 regardless of domain pair)
is explained by the fact that \textbf{all models start from the same
pre-trained initialization $\btheta_0$}.

Even though code and medical models diverge in \emph{different} directions
from $\btheta_0$, their relative displacement from the pre-trained point is
bounded by the total training budget (400 steps).  The cross-domain barrier
measures the barrier between two different ``spokes'' radiating from a
common hub ($\btheta_0$).  If the pre-trained basin is broad enough that
all 400-step trajectories remain within it, the cross-domain barrier is
just the distance across the basin, which is relatively constant.

Formally, if $\btheta_t^{\text{code}}$ and $\btheta_t^{\text{med}}$ are both within
the pre-trained basin of radius $R_0$, then the linear interpolation between
them stays within the basin, and the barrier is bounded by the basin's
internal curvature, which is domain-independent.

\section{Coexistence of Catastrophic Barriers and Good Performance}

At intermediate $\Delta W \sim 5\%$, catastrophic barriers (1.043) coexist
with best self-domain performance.  This is the signature of a
\textbf{rugged loss landscape with multiple disconnected, equally deep minima}.

\begin{claim}
The loss landscape of the fine-tuning objective contains multiple
near-equally-deep local minima that are separated by high barriers.
During the drive phase, SGD with different seeds discovers different
minima.  All discovered minima have similarly low training loss (good
self-domain performance), but the linear interpolation between minima in
different basins must cross a high-loss region (the barrier of 1.043).

This is the classic signature of \textbf{loss landscape fragmentation}
under strong label noise: the noise breaks the convexity of the basin
around $\btheta_0$, creating many sharp, disconnected minima.  The putt
phase is too short to allow barrier-crossing transitions between these
minima.
\end{claim}

\section{Quantitative Summary}

\begin{center}
\begin{tabular}{lccc}
\hline
\textbf{Quantity} & \textbf{Symbol} & \textbf{Code} & \textbf{Medical} \\
\hline
Label noise rate & $\varepsilon$ & $\lesssim 0.05$ & $\sim 0.3$--$0.5$ \\
Gradient noise inflation & $\|\bSigma_{\text{label}}\|$ & baseline & $5$--$10\times$ \\
Drive-phase $\Delta W$ & -- & $\sim 1.5\%$ & larger (est.\ $\sim 5\%$) \\
Putt-phase P\'eclet number & Pe & $\gg 1$ & $\sim 1$ \\
Barrier trajectory & $B(t)$ & inverted-U (peak 0.053 $\to$ 0.033) & monotonic (0.173 $\to$ 0.218) \\
Within-domain barrier & $\langle B_{\text{intra}}\rangle$ & 0.048 & 0.147 \\
Cross-domain barrier & $\langle B_{\text{cross}}\rangle$ & $\sim$0.05 & $\sim$0.05 \\
Early-layer barrier fraction & -- & $\sim$75\% & $\sim$75\% \\
\hline
\end{tabular}
\end{center}

\section{Testable Predictions}

\begin{enumerate}
\item \textbf{Intervention on label noise}: If we synthetically inject
      label noise into code-domain data (flip 30--50\% of labels), the
      code-domain barrier should transition from inverted-U to monotonic,
      and the within-domain barrier should rise toward the medical level
      (0.147).
\item \textbf{Extended putt phase}: Running the putt phase for $\gg 200$
      steps (e.g., 2000 steps at $\alpha = 10^{-5}$) should eventually
      cause medical barriers to decline, as even weak drift dominates
      diffusion over sufficiently long timescales.  The convergence time
      should scale as $\tau \propto 1/(\alpha \cdot \lambda_{\min})$ for
      code and $\tau \propto \mathrm{tr}(\bSigma)/(\alpha \cdot \lambda_{\min}^2)$
      for medical.
\item \textbf{Per-step gradient alignment}: Measuring the cosine similarity
      $\langle \nabla \hat{L}(\btheta_t), \nabla L(\btheta_t) \rangle$
      should reveal lower values for medical training throughout the putt
      phase, directly quantifying the gradient misalignment.
\item \textbf{Layer-wise Hessian spectrum}: The early-layer Hessian should
      exhibit larger top eigenvalues for both domains, but medical early
      layers should additionally show a \emph{wider} spectrum (more
      directions with moderate curvature), reflecting the more fragmented
      loss landscape.
\item \textbf{LMC under label smoothing}: Training medical models with
      label smoothing (soft targets instead of hard 0/1 labels) should
      reduce $\bSigma_{\text{label}}$ and produce barrier dynamics more
      similar to the code domain, potentially recovering the inverted-U
      shape.
\end{enumerate}

\section{Conclusion}

The core theoretical mechanism is:

\begin{center}
\framebox{
\begin{minipage}{0.9\textwidth}
\textbf{Label noise} $\varepsilon_{\text{med}} \gg \varepsilon_{\text{code}}$
$\implies$
\textbf{Gradient noise covariance} $\bSigma_{\text{med}} \gg \bSigma_{\text{code}}$
$\implies$
\textbf{Drive phase}: larger, less structured weight divergence
$\implies$
\textbf{Putt phase}: P\'eclet number $\mathrm{Pe}_{\text{med}} \ll \mathrm{Pe}_{\text{code}}$
$\implies$
Diffusion competes with drift $\implies$ \textbf{no basin reconnection}
$\implies$
\textbf{Monotonic barrier} (medical) vs.\ \textbf{inverted-U barrier} (code).
\end{minipage}
}
\end{center}

The model unifies all seven empirical phenomena under a single mechanism
(noise-inflated gradient covariance) and makes five falsifiable predictions.

\end{document}
